The wise virgins fall asleep. Jesus tells us that plainly. The bridegroom is delayed, the night wears on, and all ten sleep. When the cry announces his arrival, all ten wake. Five go into the wedding with him. Their sleep has not prevented them from being ready.

That detail matters when we hear Jesus tell us to watch. We might understand watchfulness as constant strain, as though faithfulness meant never becoming tired or never letting our attention rest. But the wise have rested too. The difference appears earlier in the story. They have brought oil in vessels with their lamps. They have prepared for a wait whose length they cannot control.

Saint Augustine follows that difference closely. He understands the lamps as good works that people can see. Among the examples he gives are feeding a hungry person and welcoming a stranger. Even the foolish can have these lamps. Their failure is therefore more searching than the absence of anything good in their lives. A hungry person can receive real help from someone who gives chiefly to be admired. The need has been met, while the giver's heart still needs conversion.

Augustine understands the oil as charity: the love of God and neighbor from which our actions should proceed. The vessels matter because the wise bring the oil with them. Love must truly be present within the person who acts. Other people may see what we do and thank us for it. They cannot give us the love with which we should have done it.

This is a difficult teaching for people who are trying to live faithfully. It asks us to examine something beyond whether our conduct can be defended before others. We may have done what was required and done it well. Yet the thought of receiving no recognition can trouble us more than the thought that the person we helped remains in need. The good act then reveals another part of our life that needs God's mercy.

We should be careful about what follows from discovering that desire for praise. The person we serve still needs our help. To withdraw until our motives are perfectly pure would leave that person bearing the cost of our self-examination. The work should continue as we ask God to deepen our love. We can become more attentive to the person's good, and less occupied with what the service says about ourselves. Conversion reaches into the purpose of an action that is already worth doing.

The request for oil becomes clearer in this light. The foolish turn to their companions when their lamps are failing. Their companions cannot provide what is missing. According to Augustine, the foolish have relied on other people's approval, and they seek that support again at the moment of judgment. But a companion's favorable account cannot create charity within another heart. The bridegroom's arrival discloses what each person has brought.

The refusal in the story can sound harsh. Yet it concerns the final encounter, when each person answers for a life already lived. It gives us no permission to refuse someone who needs help now. While we can still help, the need remains a claim upon us. The warning addresses our own conversion, which cannot be postponed indefinitely on the assumption that another opportunity will come.

Nor does Augustine make the wise self-satisfied. Their concern about having enough oil expresses humility before God. Even a person who has acted from love still needs forgiveness. God knows us more deeply than we know ourselves. A clear conscience can give us confidence to approach him, while we continue to depend upon his mercy. The Gospel asks for truth in our hearts without asking us to pronounce our own final verdict.

The first reading gives us a reason to make that approach with hope. Wisdom goes out to meet those who seek her. The seeker rises early and finds her already near. There is real effort in seeking, but Wisdom's approach comes before the seeker can claim to have found her. We begin to respond to God because his generosity has already reached us.

Augustine makes the same point about the oil. We can put oil into a vessel, but we did not create the olive from which it comes. Charity is God's gift. When we discover a lack of love within our service, we can ask for what we lack. We are being invited to receive and use a gift that must become real in us. Our failures give us something definite to bring before God.

The opening prayer of this Mass asks for freedom to serve him. That petition includes the whole person. Some obstacles burden the body; others affect what we desire. Dependence on admiration can make us unwilling to serve when no one will notice. As we ask God to free us for his service, we can include this impediment too. We need his help to care for someone whose gratitude may never come.

At the preparation of the gifts, the Church prays about the Passion of Christ celebrated sacramentally and received inwardly. Christ's self-giving concerns the heart with which we receive him and the life we live afterward. We approach the Eucharist as people who need his gift. The desire to be admired may still accompany us. We can acknowledge it before him and receive his help toward a love that accepts a cost for another person's good.

This need for help continues after Communion. The Church gives thanks and asks the Lord to sustain, through his Spirit, the integrity of those who have received. That continuing petition is fitting. A good intention today does not remove tomorrow's temptation, and one generous action does not finish our conversion. We depend on the Spirit as we continue the service we have begun.

Paul also speaks today to people who need hope. Some members of the Christian community have died, and those who remain are grieving. He begins with Jesus' death and resurrection. Their hope for the dead rests on what has happened to Jesus. Death has not placed the faithful beyond his care.

Grief can remain within that hope. Saint Thomas Aquinas explains that death brings the pain of bodily separation. We miss someone who is no longer beside us. Belief in the resurrection does not require us to conceal that loss. The Lord who rose from the dead gives us a future for which we can hope even while the separation hurts.

Augustine hears Paul's language about the dead in the sleep of the ten virgins. The wise also die. Their entrance into the wedding depends on a charity that has endured, rather than on their escaping death. Jesus' warning is therefore spoken by the Lord who can bring his people through death to himself. We prepare to meet someone who has already given his life for us.

You can respond to this Gospel in the service you are already giving. Where you have begun to resent a lack of thanks, continue the help that is needed, asking God to make the other person's good your purpose. That is a place where the love Augustine describes can grow. You need not wait until you feel beyond temptation to begin.

The wedding in Jesus' parable is a promised welcome as well as a judgment. The wise enter with the bridegroom. Our hope is to be received by Christ, who knows the whole truth of our lives. He gives the charity for which he asks, and his Spirit sustains us as we learn to love. We can entrust ourselves to him now, with the service still before us and the forgiveness we still need, in hope of sharing his life beyond death.
